\documentclass[hidelinks,a4paper]{jlreq}
% \def\Cjascale{0.924715}
% \def\Cjascale{0.903375}
\usepackage{luatexja}
\usepackage[haranoaji]{luatexja-preset}
% \usepackage[margin=3cm]{geometry}
\usepackage{hyperref}
\usepackage{natbib}
\bibpunct{(}{)}{;}{}{}{,}
\usepackage{xcolor}
\definecolor{hokudaiGreen}{cmyk}{0.8, 0.1, 1, 0.3}
\hypersetup{
  pdftitle={政治学入門（国際関係）参考文献},
  pdfauthor={土井翔平},
  colorlinks=true,
  setpagesize=false,
  % allcolors = blue
  allcolors = hokudaiGreen
  % linkcolor=blue,
  % filecolor=blue,      
  % urlcolor=cyan,
  % citecolor=teal,
}

\title{
    国際関係論に関する参考文献
}

\date{\today}

\author{
  土井~翔平
  \thanks{
    北海道大学大学院公共政策学連携研究部 准教授. 
    E-mail: \href{mailto:sdoi@juris.hokudai.ac.jp}{sdoi@juris.hokudai.ac.jp}
  }
}

\begin{document}

\maketitle

\section{全般}

この授業では特定の教科書は用いないが、自学自習のためにいくつかの参考書を紹介する。
ただし、国際関係論のテキストは多種多様なので、ぜひここで挙げているもの以外にもいろいろと読み比べてほしい。
いきなり一冊の教科書を読むのは面倒という場合は、政治学の教科書（例えば\citet{sunahara2020}や\citet{sakamoto2020}など）の中にある国際関係の章を読むと良い。
入門レベルの教科書としては\citet{murata2023}、中級レベルとしては\citet{kusano2023}や\citet{tago2024}などが挙げられるが、これらに限られるものではない。
より高度な教科書としては\citet{nakanishi2013}や\citet{yamakage2012}などもある。
また国際政治の重要事項に関して辞書的に利用できるものとして\citet{tanaka2010}がある（ただし、やや古い点に注意）。
国際関係論に関して卒業論文を執筆したい、あるいは大学院に進学したいと考えている場合は英語の教科書として\citet{frieden2019}や\citet{baylis2026}などをおすすめする。

また、国際関係を理解するためには国際関係史、国際法、比較政治学といった隣接分野の知識も必要である。
これらについても様々なテキストがあるが、国際政治史であれば\citet{ogawa2018}が定評のある入門書である。
国際法の入門書としては\citet{tamada2022}があり、\citet{kitamura2025}は身近なテーマから国際法について平易に説明してくれる。
比較政治学の入門書としては、民主主義国について\citet{ito2025}、非民主主義国について\citet{kubo2025}がある。
この他にも国際関係論や政治学を理解する上でゲーム理論や計量分析の知識が有用になる局面はあるが、これらは後述する。

\section{国際社会}

\section{安全保障}

\section{国際政治経済}

\section{ゲーム理論}

\section{計量分析}

\bibliography{references}
\bibliographystyle{jecon}

\end{document}